\input{estilo}
% No sigue el orden de la pizarra: el two-phase commit de Spanner como lo
% describe el paper (OSDI 2012, §4.2.1), sin los timestamps, que el apunte
% presenta recién en las secciones 5 y 6. Las lecturas en tinta, el commit en
% verde y los locks en rojo, como en la pizarra.
\begin{document}
\begin{figura}\begin{tikzpicture}[y=0.9cm]
  \def\C{0} \def\Pa{4.8} \def\Pb{9.6}
  % cabeceras: el cliente y los líderes de dos grupos de Paxos
  \node at (\C,0.35) {\var{C}};
  \foreach \x/\n in {\Pa/1, \Pb/2} {
    \foreach \d in {-1,0,1} \draw[fill=relleno] (\x+\d*0.42-0.14,1.05) rectangle ++(0.28,0.28);
    \fill[acento] (\x,1.19) circle (1.6pt);
    \draw[llave=4pt] (\x+0.6,0.92) -- (\x-0.6,0.92);
    \node at (\x,0.35) {\var{P}\textsubscript{\n}};
  }
  \foreach \x in {\C,\Pa,\Pb} \draw[vida] (\x,0) -- (\x,-7.9);
  \tikzset{lock/.style={etiq, text=acento, anchor=west}, nota/.style={etiq, align=left, anchor=west}}
  % las lecturas, con read locks en el líder
  \draw[msg] (\C,-0.6) -- node[etiq, above] {read \var{x}} (\Pa,-0.6);
  \node[evento=acento] at (\Pa,-0.6) {}; \node[lock] at (\Pa+0.1,-0.6) {read lock};
  \draw[respuesta] (\Pa,-0.75) -- node[etiq, below, sloped] {\var{x}} (\C,-1.1);
  \draw[msg] (\C,-1.5) -- node[etiq, above, pos=0.7] {read \var{y}} (\Pb,-1.5);
  \node[evento=acento] at (\Pb,-1.5) {}; \node[lock] at (\Pb+0.1,-1.5) {read lock};
  \draw[respuesta] (\Pb,-1.65) -- node[etiq, below, sloped, pos=0.3] {\var{y}} (\C,-2.1);
  \node[evento] at (\C,-2.55) {};
  \node[etiq, align=right, anchor=east] at (\C-0.15,-2.55) {calcula y guarda\\los writes en un buffer};
  % separador de las dos fases
  \draw[suave, line width=0.4pt, dash pattern=on 1pt off 2pt] (-3.2,-3.0) -- (13.6,-3.0);
  \node[etiq, text=suave, anchor=south east] at (13.6,-3.0) {lecturas};
  \node[etiq, text=suave, anchor=north east] at (13.6,-3.0) {two-phase commit};
  % el cliente manda el commit, con los writes, a cada participante
  \begin{scope}[every node/.style={etiq, text=segundo}]
    \draw[msg, segundo] (\C,-3.4) -- node[above, pos=0.72, sloped] {commit + writes, TC = \var{P}\textsubscript{1}} (\Pb,-3.75);
    \draw[msg, segundo] (\C,-3.6) -- node[below, sloped] {commit + writes, TC = \var{P}\textsubscript{1}} (\Pa,-3.85);
  \end{scope}
  \node[evento=acento] at (\Pb,-3.75) {}; \node[lock] at (\Pb+0.1,-3.75) {write lock};
  \node[evento=acento] at (\Pa,-3.85) {}; \node[lock] at (\Pa+0.1,-4.0) {write lock};
  \begin{scope}[every node/.style={nota, text=segundo}]
    % el participante se prepara solo
    \node[evento=segundo, anchor=center] at (\Pb,-4.45) {};
    \node at (\Pb+0.1,-4.55) {prepare:\\lo registra en Paxos};
    \draw[respuesta, segundo] (\Pb,-4.95) -- node[below, sloped, anchor=north] {preparado} (\Pa,-5.3);
    \node[etiq, text=suave, anchor=east] at (\Pa-0.1,-4.45) {el TC no hace prepare};
    % el coordinador decide
    \node[evento=segundo, anchor=center] at (\Pa,-5.45) {};
    \node[anchor=east, align=right] at (\Pa-0.15,-5.55) {con todos preparados,\\registra el commit en Paxos};
    % y avisa a todos
    \draw[msg, segundo] (\Pa,-6.1) -- node[below, sloped] {commit} (\C,-6.45);
    \draw[msg, segundo] (\Pa,-6.1) -- node[above, sloped] {commit} (\Pb,-6.45);
    \node[evento=segundo, anchor=center] at (\Pb,-6.45) {};
    \node at (\Pb+0.1,-6.45) {registra el commit en Paxos};
    \node[evento=segundo, anchor=center] at (\Pa,-7.3) {}; \node[evento=segundo, anchor=center] at (\Pb,-7.3) {};
    \node at (\Pa+0.1,-7.3) {aplica y libera locks};
    \node (lib) at (\Pb+0.1,-7.3) {aplica y libera locks};
  \end{scope}
  \draw[<-] ($(lib.east)+(0.05,0)$) -- ++(0.5,0) node[etiq, right] {2PL};
\end{tikzpicture}\end{figura}
\end{document}
